% ECONOMICS_PROBLEM: {"id": "D81-567", "title": "Risk preferences versus complexity errors", "jel_code": "D81", "source_id": "OP-0300"}
% !TeX program = xelatex
\documentclass[11pt,a4paper]{article}
\usepackage[margin=23mm]{geometry}
\usepackage{fontspec}
\setmainfont{Latin Modern Roman}
\usepackage{amsmath,amssymb,mathtools}
\usepackage{xcolor,enumitem,booktabs,longtable,array,xurl,hyperref}
\definecolor{muted}{HTML}{53636C}
\hypersetup{colorlinks=true,urlcolor=blue,linkcolor=blue}
\setlength{\parindent}{0pt}
\setlength{\parskip}{5pt}
\newcommand{\R}{\mathbb R}
\newcommand{\E}{\mathbb E}
\newcommand{\N}{\mathbb N}
\newcommand{\ind}{\mathbf 1}
\newcommand{\Var}{\operatorname{Var}}
\newcommand{\Cov}{\operatorname{Cov}}
\newcommand{\supp}{\operatorname{supp}}
\DeclareMathOperator*{\argmin}{arg\,min}
\DeclareMathOperator*{\argmax}{arg\,max}
\newcommand{\entrymeta}[3]{{\sffamily\small\color{muted}\textbf{\texttt{#1}}\quad|\quad #2\par #3\par}}
\newcommand{\Problemref}[1]{\texttt{#1}}
\newcommand{\JELref}[2]{\texttt{#2} (JEL \texttt{#1})}
\begin{document}
\section*{Risk preferences versus complexity errors}
\textbf{Problem ID:} \texttt{D81-567}\quad \textbf{JEL:} \texttt{D81}\par
% BEGIN ECONOMICS STATEMENT
\label{problem:OP-0300}
{\sffamily\small\color{muted}Original subject group: Behavioral and experimental economics\par}
\label{definitions:problem:30007582}
\textbf{Audit classification:} Published research agenda. Full Oprea 2024 article now retrieved. §§III–IV, pp.3806–3807, explicitly proposes further investigation of heuristic behavior and risk-induced complexity. The identification/field decomposition is editorial.
\subsection*{Definitions and precise research target}
\textbf{Model and research target.} Consider individuals \(i\) choosing between a sure payment and a binary lottery \(L=(p,x,y)\), where \(p\in[0,1]\) is the probability of receiving \(x\), and \(y\) is received otherwise. Let \(C\in\{0,1\}\) randomly assign a simple or complex representation while preserving \(L\), payment rules, and information. Let \(a_i(L,C)\) be the reported certainty equivalent. A preference model specifies utility \(u_i\), probability weighting \(w_i\), and a reference point \(r_i\); its latent certainty equivalent \(v_i(L)\) satisfies \(u_i(v_i)=w_i(p)u_i(x)+(1-w_i(p))u_i(y)\). Observations obey \(a_i(L,C)=v_i(L)+e_i(L,C)\), where \(e_i\) is a decision or reporting error whose conditional mean need not vanish. The observable causal complexity effect is
\[\tau(L)=\mathbb E[a_i(L,1)-a_i(L,0)].\]
Identify which restrictions, auxiliary tasks, or repeated measurements permit separation of \(v_i\) from systematic \(e_i\), rather than attributing every pattern to preferences. Random assignment identifies \(\tau\) but does not identify the error-free preference distribution without further assumptions. A field extension should specify a target population \(P\), comparable lotteries, and consequential implementation; estimate the corresponding \(\tau_P(L)\) and the predictive performance of recovered preferences for held-out decisions.

\emph{Definitions and maintained assumptions (editorial unless a source definition is named).} For a binary gain lottery with \(x\ge y\), strictly increasing utility \(u_i\) and weighting \(w_i:[0,1]\to[0,1]\), \(w_i(0)=0,w_i(1)=1\), the specified model defines a unique certainty equivalent \(v_i(L)\in[y,x]\) by the displayed equation. A reference-dependent version must explicitly replace \(u_i(x)\) by \(u_i(x-r_i)\) and specify treatment-invariance of \(r_i\); merely naming a reference point does not use it. Measurements are either monetary certainty equivalents or choices converted by a declared elicitation procedure and error kernel. Potential responses \(a_i(L,C)\) exist for both equivalent representations, with randomized \(C\). To infer preferences, assume representations preserve utility, beliefs and reference points; otherwise \(\tau(L)\) is a representation effect with several channels. Nonzero systematic error permits arbitrary offsets between \(v_i\) and \(e_i\), so point identification needs anchors, repeated-treatment restrictions or a parametrized error law. A deterministic mirror task is a separate economic task and its validity as an error proxy must be tested.
\subsection*{Variants and assumption changes}
\begin{enumerate}
\item \textbf{Repeated representation (editorial).} Repeatedly elicit identical lotteries under randomized equivalent descriptions, with calibrated payment procedures. Identify a maintained stable utility/weighting model under explicit error restrictions, and report failure when systematic errors are unrestricted.
\item \textbf{Risk and mirrors (source-related editorial).} Jointly vary risk presence and representation complexity using matched tasks. Test whether mirror errors predict lottery errors and bound decomposition sensitivity to risk-dependent error or reference shifts.
\end{enumerate}
\subsection*{References and source audit}
\begin{enumerate}
\item §§III–IV pp.3806–3807 in published article. Supports the stated research area; exact displayed specification is editorial. \url{https://www.aeaweb.org/articles?id=10.1257/aer.20221227}
\item §§III–IV pp.3806–3807 in published article. Supports the stated research area; exact displayed specification is editorial. \url{https://ryanoprea.com/papers/decisions-under-risk.pdf}
\end{enumerate}
\begin{enumerate}\raggedright
\item\hypertarget{definitions:source:30007582:1}{} Ryan Oprea. 2024. \emph{Decisions under Risk Are Decisions under Complexity}. §§III–IV, AER114(12), pp.3806–3807; full article retrieved from author: https://ryanoprea.com/papers/decisions-under-risk.pdf.
\url{https://www.aeaweb.org/articles?id=10.1257/aer.20221227}.
DOI: \url{https://doi.org/10.1257/aer.20221227}.
\item\hypertarget{definitions:source:30007582:2}{} Ryan Oprea. 2024. \emph{Decisions under Risk Are Decisions under Complexity}. §§III–IV, pp.3806–3807.
\url{https://ryanoprea.com/papers/decisions-under-risk.pdf}.
DOI: \url{https://doi.org/10.1257/aer.20221227}.
\end{enumerate}

\subsection*{Common conventions: Source mathematical conventions}

These conventions apply to each entry unless explicitly replaced.

\textbf{Models and identification.} A primitive model \(m\) specifies all probability laws, feasible choices, information, equilibrium and measurement rules used by its target. The admissible class \(\mathcal M\) is fixed independently of outcomes. If \(P_O(m)\) is its observable probability law and \(\tau(m)\) the target, its identified set at observable law \(P\) is
\[
 \mathcal I_\tau(P)=\{\tau(m):m\in\mathcal M,\ P_O(m)=P\}.
\]
Point identification means that this set is a singleton. An empty set signals incompatibility of data and assumptions. Coordinate bounds need not describe the joint set. Bounds are sharp only if every retained target is attainable or arbitrarily approachable by compatible models. Finite bounds require explicit restrictions on unobserved quantities; unrestricted confounding, hidden wealth, extrapolation or arbitrary equilibrium selection can yield uninformative sets.

\textbf{Causal interventions.} A potential outcome \(Y_i(p)\) is the outcome under a complete policy regime \(p\), including timing, financing, information and market spillovers. The target population and units are fixed before treatment. With interference, use the complete assignment vector or a justified exposure mapping; individual no-interference notation alone cannot identify an economy-wide intervention. A difference of mean potential outcomes does not require the cross-world joint distribution. Conditioning on post-treatment migration, survival, enrollment or employment changes the estimand and needs separate justification.

\textbf{Statistical statements.} An estimator, confidence set, forecast or test specifies a sampling law, model class and loss/coverage criterion. Uniform \((1-\alpha)\)-coverage means \(\inf_{m\in\mathcal M}\Pr_m\{\tau(m)\in C_n\}\geq1-\alpha\), or an explicitly asymptotic version. The whole parameter space trivially has coverage, so informativeness requires a stated width, risk or power objective. A forecasting improvement is relative to a named benchmark, information set and evaluation population.

\textbf{Equilibrium and optimization.} An equilibrium comprises feasible plans, optimality, beliefs where needed, budget constraints and all market/resource clearing. The primitives or a stated benchmark define each correspondence \(\mathcal E(p;m)\). Nonemptiness is assumed or proved, never inferred just from compact policy sets. A selected-equilibrium target uses a declared selection rule; a robust target quantifies over all permitted equilibria. Use suprema or infima unless attainment is established. An \(\varepsilon\)-optimal policy is feasible and within \(\varepsilon\) of the finite optimum value. Report an empty feasible set separately.

\textbf{Welfare and accounting.} Normative utility, interpersonal normalization, social weights, numeraire, discounting and the use of public receipts are primitives. Behavioral utility need not coincide with the welfare criterion. Household welfare includes ownership income and rents once; adding profits again to compensating variation generally double-counts them. A monetary equivalent is defined by an explicit transfer and price convention, with monotonicity and range conditions. Average effects, mechanism contributions, transfers and welfare losses are different targets. Infinite sums require convergence and dynamic feasible plans require terminal or no-Ponzi conditions.

\textbf{Source versions.} A working-paper date, revision date and journal publication year are distinguished. A DOI for a journal version is not automatically the DOI of an earlier manuscript. Locators refer to the stated version and printed pages where specified. A metadata-only check does not verify a claim about a full-text conclusion. Every entry's source subsection narrows the provenance claim accordingly.
% END ECONOMICS STATEMENT
\end{document}
